% =====================================================================
%  Appendix module: Posterior distribution of RMSE
%  Modular subfile: compiles standalone (uses the main preamble) OR is
%  pulled into the main document via \subfile{...}. See main file header.
%  Figure built by: manuscript/figures/build_appendix_rmse_posterior.R
%  Numbers also written to: manuscript/figures/appendix_rmse_posterior.csv
% =====================================================================
\documentclass[../HIKET_main_manuscript.tex]{subfiles}

\begin{document}

\section{Posterior distribution of RMSE}
\label{app:rmse-posterior}

The multimodel metrics table reports a single RMSE per model, computed from the
posterior-mean prediction. That number says nothing about how tightly the fit is
pinned: two calibrations with the same central RMSE can disagree with themselves
by very different amounts. Here RMSE is recomputed for \emph{every} posterior
draw, so each panel of Figure~\ref{fig:rmse-posterior} is a distribution whose
location is the fit and whose width is the posterior's own disagreement about it.
Calibration and held-out plots are drawn together, because the gap between them
is the part that is not merely a fit statistic.

\begin{figure}[htbp]
  \centering
  \includegraphics[width=\linewidth]{figures/appendix_rmse_posterior.png}
  \caption{Posterior distribution of RMSE, one panel per model, recomputed for each
  of the 100 posterior draws. Blue = calibration plots, red = held-out plots;
  dashed lines are medians; densities are scaled to a common height so shape is
  comparable across panels. The horizontal range is held fixed across runs so that
  successive calibrations can be laid side by side.}
  \label{fig:rmse-posterior}
\end{figure}

Two features are worth noting. First, the six models are barely separated in
central RMSE (calibration medians span $43.6$--$44.4$~tC\,ha$^{-1}$), which is the
same message the skill metrics carry: structure buys very little here. Second, and
less visible in a point estimate, the calibration--holdout gap is \emph{not}
uniform across the ensemble. TP2 and TP3 achieve the best calibration RMSE of the
six yet the worst on held-out plots ($43.6 \rightarrow 46.8$ and
$43.6 \rightarrow 46.9$), while Yasso15's two distributions almost coincide
($44.2 \rightarrow 44.4$). The models that fit the calibration set best are
therefore the ones that generalise worst, and they are also the models carrying
the largest input multipliers --- consistent with fit being bought by rescaling
inputs rather than by resolving kinetics.

This figure is intended as a comparison instrument across calibrations. Where a
prior is tightened onto a physical referent, the expectation is that fit degrades;
the quantity of interest is then not only how far the distribution moves but
whether its width moves with it, since a constraint that shifts the location while
leaving the spread untouched is doing something different from one that does both.

\end{document}
